% Minimal preamble for Open Logic Project sources (pinned commit 1e960beff9ed7835bf3e3f1335e21af3439cd107).
%
% OLP's own master preamble cannot run on the project host: TeX Live 2025 there lacks
% prooftree.sty, ntheorem.sty, extarrows.sty, mathdots.sty, isomath.sty and subfiles.sty.
% Following the project's own rule ("if the author preamble needs packages the toolchain lacks,
% do not invent macros -- use a minimal working preamble"), this file declares, verbatim in
% spirit and numbering behaviour, exactly the theorem-like environments OLP defines in
% sty/open-logic-book-envs.sty (thm/lem/prop/cor/defn/ex/prob/rem/note/case/axiom/conv),
% plus the shorthand accents OLP uses (`!!a`, `!!b`, `!!{...}`).
%
% Author text is untouched; only the packaging is ours.
\documentclass[10pt,a4paper]{article}

\usepackage{amsmath}
\usepackage{amssymb}
\usepackage{amsthm}
\usepackage{graphicx}
\usepackage{hyperref}
\usepackage{enumitem}
\usepackage[framemethod=default]{mdframed}

% ---- OLP shorthand accents (from sty/open-logic.sty) ----
% NB: \! is taken by LaTeX (negative thin space), so we only define the !!a / !!b families
% and the \olup helper used by them.
\newcommand*{\olup}[1]{\MakeUppercase{#1}}
\def\!!!a#1{\olup{#1}}
\def\!a#1{\olup{#1}}
\def\!b#1{\olup{#1}}

% ---- theorem-like environments, same family numbering as OLP ----
\theoremstyle{plain}
\newtheorem{thm}{Theorem}[section]
\newtheorem{lem}[thm]{Lemma}
\newtheorem{prop}[thm]{Proposition}
\newtheorem{cor}[thm]{Corollary}
\newtheorem{axiom}{Axiom}

\theoremstyle{definition}
\newtheorem{defn}[thm]{Definition}
\newtheorem{ex}[thm]{Example}
\newtheorem{prob}{Problem}[section]
\newtheorem{probd}{Problem}
\newtheorem{conv}[thm]{Convention}

\theoremstyle{remark}
\newtheorem{rem}[thm]{Remark}
\newtheorem{note}[thm]{Note}
\newtheorem{case}[thm]{Case}

% ---- environments OLP defines as wrappers ----
\newenvironment{probtag}[1]{\begin{prob}}{\end{prob}}
\newenvironment{explain}{}{}
\newenvironment{explain*}{}{}
\newenvironment{editorial}{}{}
\newenvironment{digress}{}{}
\newenvironment{digress*}{}{}
\newenvironment{history}{}{}
\newenvironment{reference}{}{}
\newcommand*{\olfileid}[3]{}
\newcommand*{\olsection}[1]{\section*{#1}}
\newcommand*{\olchapter}[3]{}
\newcommand*{\olpart}[2]{}
\newcommand*{\olimport}[1]{}
\newcommand*{\olinclude}[1]{}
\newcommand*{\olpath}[1]{}
\newcommand*{\ollangid}[1]{}
\newcommand*{\problemtag}[1]{}
\newcommand*{\problemauthor}[1]{}
\newcommand*{\ollabel}[1]{\label{#1}}
% \olref takes 0-3 OPTIONAL arguments plus one mandatory label:
%     \olref{lab}  |  \olref[from]{lab}  |  \olref[from][to]{lab}  |  \olref[a][b][c]{lab}
% A single-argument definition silently swallows "[sth]" and emits a broken reference
% (measured on the host: the text came out as "??iso]x"). The chain below consumes the
% brackets and cross-references only the mandatory argument.
\makeatletter
\newcommand{\olref}{\@ifnextchar[{\olref@i}{\olref@final}}
\def\olref@i[#1]{\@ifnextchar[{\olref@ii}{\olref@final}}
\def\olref@ii[#1]{\@ifnextchar[{\olref@iii}{\olref@final}}
\def\olref@iii[#1]{\olref@final}
\def\olref@final#1{\ref{#1}}
\makeatother
\newcommand*{\olnewpage}{\clearpage}

% ===================== OLP compatibility layer (1) =====================

\makeatletter

% ---------------- proof trees (OLP's prooftree is a bussproofs derivative) ----------------
\providecommand{\AxiomC}[1]{}
\providecommand{\UnaryInfC}[1]{}
\providecommand{\BinaryInfC}[1]{}
\providecommand{\TrinaryInfC}[1]{}
\providecommand{\QuaternaryInfC}[1]{}
\providecommand{\QuinaryInfC}[1]{}
\providecommand{\RightLabel}[1]{}
\providecommand{\LeftLabel}[1]{}
\providecommand{\doubleLine}{}
\providecommand{\noLine}{}
\providecommand{\DisplayProof}{}
\providecommand{\fCenter}{}
\newenvironment{prooftree}{}{}
\newenvironment{ltprooftree}{}{}
\newenvironment{oltableau}{}{}

% ---------------- logical notation ----------------
\providecommand{\lif}{\mathbin{\rightarrow}}
\providecommand{\lforall}{\mathord{\forall}}
\providecommand{\lexists}{\mathord{\exists}}
\providecommand{\ltrue}{\top}
\providecommand{\lfalse}{\bot}
\providecommand{\Ltrue}{\top}
\providecommand{\Lfalse}{\bot}
\providecommand{\Sequent}{\mathrel{\Rightarrow}}
\providecommand{\Proves}{\mathrel{\vdash}}
\providecommand{\Entails}{\mathrel{\models}}
\providecommand{\Subst}[3]{#1[#2/#3]}
\providecommand{\sFmla}[2]{#1(#2)}
\providecommand{\pFmla}[2]{#1(#2)}
\providecommand{\formula}{\relax}
\providecommand{\True}{\top}
\providecommand{\False}{\bot}
\providecommand{\Args}[1]{#1}
\providecommand{\Ts}{\ensuremath{\top}}
\providecommand{\Fi}{\ensuremath{\bot}}
\providecommand{\lfalsum}{\bot}

% ---------------- objects, structures, metatheory ----------------
\providecommand{\Obj}{}
\providecommand{\Struct}{\ensuremath{\mathfrak{A}}}
\providecommand{\Domain}{\ensuremath{|\mathfrak{A}|}}
\providecommand{\Assign}{\ensuremath{s}}
\providecommand{\Sat}[2]{\ensuremath{#1 \models #2}}
\providecommand{\mSat}[3]{\ensuremath{#1 \models #2[#3]}}
\providecommand{\mModel}[2]{\ensuremath{#1 \models #2}}
\providecommand{\Lang}{\ensuremath{\mathcal{L}}}
\providecommand{\Th}{\ensuremath{\mathrm{Th}}}
\providecommand{\Log}{\ensuremath{\mathrm{L}}}
\providecommand{\tuple}[1]{\langle #1 \rangle}
\providecommand{\num}[1]{\underline{#1}}
\providecommand{\cfind}[1]{\varphi_{#1}}
\providecommand{\Nat}{\ensuremath{\mathbb{N}}}
\providecommand{\Mult}{\ensuremath{\mathrm{Mult}}}
\providecommand{\Add}{\ensuremath{\mathrm{Add}}}
\providecommand{\Exp}{\ensuremath{\mathrm{Exp}}}
\providecommand{\Zero}{\ensuremath{0}}
\providecommand{\Succ}{\ensuremath{\mathrm{S}}}
\providecommand{\lambd}{\ensuremath{\lambda}}
\providecommand{\bredpar}{\mathrel{\rightarrow_\beta}}
\providecommand{\bcd}[1]{\ensuremath{\ulcorner #1 \urcorner}}
\providecommand{\Set}[1]{\ensuremath{\{#1\}}}
\providecommand{\indcase}[1]{\par\noindent\emph{Case #1}.\quad}
\providecommand{\indfrm}[1]{\par\noindent\emph{Form #1}.\quad}
\providecommand{\iftag}[1]{#1}
\providecommand{\tagitem}[1]{\item[#1]}
\providecommand{\Elim}[1]{}
\providecommand{\Intro}[1]{}
\providecommand{\LeftR}[1]{}
\providecommand{\RightR}[1]{}
\providecommand{\TAss}[1]{}
\providecommand{\TRule}[1]{}

\makeatother
\providecommand{\olpath}{.}

% ===================== OLP compatibility layer (2) =====================
% --- second compatibility block: notation and shorthands used by the frozen OLP units ---
\providecommand{\CutCS}{\mathrm{Cut}}
\providecommand{\LeftSubproofLabel}{\relax}
\providecommand{\RightSubproofLabel}{\relax}
\providecommand{\maeh}{\relax}
\providecommand{\cheight}{\relax}
\providecommand{\ande}{\mathrm{and}_E}
\providecommand{\andi}{\mathrm{and}_I}
\providecommand{\ore}{\mathrm{or}_E}
\providecommand{\ori}{\mathrm{or}_I}
\providecommand{\typeof}{\mathord{:}}
\providecommand{\pValue}[1]{#1}
\providecommand{\Sent}{\mathrm{Sent}}
\providecommand{\cardplus}{\mathbin{+}}
\providecommand{\cardtimes}{\mathbin{\cdot}}
\providecommand{\cardexpo}{\mathbin{\hat{\ }}}
\providecommand{\cardnless}{\not<}
\providecommand{\cardle}{\leq}
\providecommand{\funfromto}[3]{#1\colon #2 \to #3}
\providecommand{\funimage}[1]{\operatorname{im}(#1)}
\providecommand{\fdefined}{\mathord{\downarrow}}
\providecommand{\canonord}{\operatorname{ord}}
\providecommand{\ordtimes}{\mathbin{\cdot}}
\providecommand{\ordeq}[2]{#1 = #2}
\providecommand{\ordneq}[2]{#1 \neq #2}
\providecommand{\ordtype}[1]{\operatorname{ord}(#1)}
\providecommand{\ordsucc}{\mathord{+1}}
\providecommand{\ordunion}{\cup}
\providecommand{\onesphere}{S^1}
\providecommand{\rotationsgroup}{\mathrm{SO}(2)}
\providecommand{\equivrep}[1]{[#1]}
\providecommand{\rlexless}{\prec}
\providecommand{\ZFC}{\mathrm{ZFC}}
\providecommand{\ZF}{\mathrm{ZF}}
\providecommand{\Rat}{\mathbb{Q}}
\providecommand{\Real}{\mathbb{R}}
\providecommand{\PosInt}{\mathbb{Z}^{+}}
\providecommand{\Id}{\mathrm{id}}
\providecommand{\Bin}{\mathrm{Bin}}
\providecommand{\pto}{\rightharpoonup}
\providecommand{\dom}{\operatorname{dom}}
\providecommand{\ran}{\operatorname{ran}}
\providecommand{\len}{\operatorname{len}}
\providecommand{\concat}{\mathbin{{}^\frown}}
\providecommand{\card}[1]{\lvert #1 \rvert}
\providecommand{\Pow}{\mathcal{P}}
\providecommand{\Setabs}[2]{\{\, #1 \mid #2 \,\}}
\providecommand{\comp}{\circ}
\providecommand{\Char}{\chi}
\providecommand{\Prf}{\operatorname{Prf}}
\providecommand{\Prov}{\operatorname{Prov}}
\providecommand{\Proj}{\operatorname{proj}}
\providecommand{\Gn}{\operatorname{gn}}
\providecommand{\gn}{\operatorname{gn}}
\providecommand{\Value}[1]{\operatorname{val}(#1)}
\providecommand{\Deduce}{\mathrm{D}}
\providecommand{\Dual}{\mathrm{d}}
\providecommand{\Ax}{\mathrm{Ax}}
\providecommand{\pSat}[3]{#1 \models #2\,[#3]}
\providecommand{\pAssign}{\sigma}
\providecommand{\bforall}{\forall}
\providecommand{\bexists}{\exists}
\providecommand{\liff}{\mathbin{\leftrightarrow}}
\providecommand{\FV}{\operatorname{FV}}
\providecommand{\UnaryInf}{\ensuremath{\mid}}
\providecommand{\BinaryInf}{\ensuremath{\mid}}
\providecommand{\depth}{\operatorname{depth}}
\providecommand{\oliflabeldef}[3]{#1}
\providecommand{\TMblank}{\sqcup}
\providecommand{\TMstroke}{\mid}
\providecommand{\TMstay}{\mathrm{S}}
\providecommand{\TMright}{\mathrm{R}}
\providecommand{\TMleft}{\mathrm{L}}
\providecommand{\TMendtape}{\triangleright}
\providecommand{\umin}{\mu}
\providecommand{\Sum}{\sum}
\providecommand{\citep}[1]{\cite{#1}}
\providecommand{\citet}[1]{\cite{#1}}
\providecommand{\citeauthor}[1]{\cite{#1}}
\providecommand{\citeyear}[1]{\cite{#1}}

% ===================== OLP compatibility layer (3) =====================
% --- third compatibility block: remaining notation seen in the frozen OLP units ---
\providecommand{\iso}{\cong}
\providecommand{\nszero}{\mathrm{0}}
\providecommand{\nssucc}{\mathrm{S}}
\providecommand{\nsplus}{\mathbin{+}}
\providecommand{\nssless}{<}
\providecommand{\nsless}{<}
\providecommand{\PL}{\mathrm{PL}}
\providecommand{\Nec}{\mathord{\Box}}
\providecommand{\RK}{\mathrm{RK}}
\providecommand{\ST}{\mathrm{ST}}
\providecommand{\SSubst}[3]{#1[#2/#3]}
\providecommand{\subst}[3]{#1[#2/#3]}
\providecommand{\tagfalse}{\mathrm{F}}
\providecommand{\cutr}{\mathrm{Cut}}
\providecommand{\cardneq}{\neq}
\providecommand{\issucc}[1]{\operatorname{succ}(#1)}
\providecommand{\tup}[1]{\langle #1 \rangle}
\providecommand{\Var}{\mathrm{Var}}
\providecommand{\Bew}{\operatorname{Bew}}
\providecommand{\Con}{\operatorname{Con}}
\providecommand{\Cons}{\operatorname{Cons}}
\providecommand{\pDeduce}{\mathrm{D}}
\providecommand{\psat}{\models}
\providecommand{\den}[1]{\llbracket #1 \rrbracket}
\providecommand{\Sem}[2]{\llbracket #1 \rrbracket_{#2}}
\providecommand{\DomainOf}[1]{\lvert #1 \rvert}
\providecommand{\Form}{\mathrm{Form}}
\providecommand{\Terms}{\mathrm{Term}}
\providecommand{\Atom}{\mathrm{Atom}}
\providecommand{\FVars}[1]{\operatorname{FV}(#1)}
\providecommand{\fv}[1]{\operatorname{fv}(#1)}
\providecommand{\tuplep}[1]{\langle #1 \rangle}
\providecommand{\restrict}[2]{#1 \upharpoonright #2}
\providecommand{\sto}{\rightarrow}
\providecommand{\FmlaSet}[1]{\mathrm{Fmla}_{#1}}
\providecommand{\Thm}{\mathrm{Thm}}
\providecommand{\Eval}{\operatorname{val}}
\providecommand{\blank}{\sqcup}
\providecommand{\blankTape}{\sqcup}
\providecommand{\Consis}{\operatorname{Con}}
\providecommand{\drs}{\mathrm{d}}
\providecommand{\VRule}{\mid}
\providecommand{\Bet}{\mathrm{B}}
\providecommand{\Seg}[1]{\mathrm{Seg}_{#1}}
\providecommand{\Enums}{\mathrm{En}}
\providecommand{\Z}{\mathbb{Z}}
\providecommand{\N}{\mathbb{N}}
\providecommand{\abs}[1]{\lvert #1 \rvert}
\providecommand{\floor}[1]{\lfloor #1 \rfloor}
\providecommand{\ceil}[1]{\lceil #1 \rceil}

% ===================== OLP compatibility layer (4) =====================
% --- fourth compatibility block: harvested from the per-candidate compile logs ---
\providecommand{\beredpar}{\mathrel{\rightarrow_\beta}}
\makeatletter
\def\bered{\@ifnextchar[{\bered@i}{\mathrel{\rightarrow_\beta}}}
\def\bered@i[#1]{\mathrel{\rightarrow_\beta}}
\def\bred{\@ifnextchar[{\bred@i}{\mathrel{\rightarrow_\beta}}}
\def\bred@i[#1]{\mathrel{\rightarrow_\beta}}
\makeatother
\providecommand{\beredone}{\mathrel{=_\beta}}
\providecommand{\bredone}{\mathrel{=_\beta}}
\providecommand{\xredone}{\mathrel{=_\beta}}
\providecommand{\aconvone}{\mathrel{=_\alpha}}
\providecommand{\cardless}{<}
\providecommand{\tagrefs}[1]{#1}
\providecommand{\ordplus}{\mathbin{+}}
\providecommand{\funrestrictionto}[2]{#1 \upharpoonright #2}
\providecommand{\fundefined}{\mathord{\uparrow}}
\providecommand{\Trm}{\mathrm{Trm}}
\providecommand{\OProv}{\operatorname{Prov}}
\providecommand{\OPrf}{\operatorname{Prf}}
\providecommand{\aeq}{\mathrel{=}}
\providecommand{\mClass}[1]{[#1]}
\providecommand{\equivclass}[1]{[#1]}
\providecommand{\pheight}{\operatorname{ht}}
\providecommand{\phy}{\operatorname{ht}}
\providecommand{\Cut}{\mathrm{Cut}}
\providecommand{\QuantRank}{\operatorname{qr}}
\providecommand{\elemequiv}{\mathrel{\equiv}}
\providecommand{\PIso}{\mathrm{PIso}}
\providecommand{\Theory}{\mathrm{Th}}
\providecommand{\supstrict}{\subsetneq}
\providecommand{\trcl}{\operatorname{trcl}}
\providecommand{\cref}[1]{\ref{#1}}
\providecommand{\Discharge}{\mathrm{Dis}}
\providecommand{\Weakening}{\mathrm{Wk}}
\providecommand{\ZFminus}{\mathrm{ZF}^{-}}
\makeatletter
\renewcommand{\Dual}{\@ifnextchar[{\Dual@i}{\ifmmode\mathrm{d}\else$\mathrm{d}$\fi}}
\def\Dual@i[#1]{\ifmmode\mathrm{d}[#1]\else$\mathrm{d}[#1]$\fi}
\makeatother
\providecommand{\RightR}{\mathrm{R}}
\providecommand{\LeftR}{\mathrm{L}}
\providecommand{\Intro}{\mathrm{I}}
\providecommand{\Axiom}{\mathrm{Ax}}
\providecommand{\formula}[1]{#1}
\providecommand{\stared}{\ast}
\providecommand{\oset}[2]{\{\, #1 \mid #2 \,\}}
\providecommand{\setof}[2]{\{\, #1 \mid #2 \,\}}
\providecommand{\pow}[1]{\mathcal{P}(#1)}
\providecommand{\ordpair}[2]{\langle #1, #2 \rangle}
\providecommand{\ordtriple}[3]{\langle #1, #2, #3 \rangle}
\providecommand{\union}{\cup}
\providecommand{\inter}{\cap}
\providecommand{\seq}[1]{\langle #1 \rangle}
\providecommand{\SuchThat}{\mid}
\providecommand{\where}{\mid}
\providecommand{\TVar}{\mathrm{TVar}}
\providecommand{\Bool}{\mathrm{Bool}}
\providecommand{\Val}{\operatorname{Val}}
\providecommand{\TL}[1]{\mathrm{TL}}
\providecommand{\xstared}{\ast}

% ===================== OLP compatibility layer (5) =====================
% --- fifth compatibility block: rescues the quarantined OLP candidates ---
\providecommand{\ident}{\operatorname{id}}
\providecommand{\setrank}{\operatorname{rank}}
\providecommand{\cardfont}{\mathfrak{c}}
\providecommand{\Complement}{\overline}
\providecommand{\QR}{\mathrm{QR}}
\providecommand{\becd}{\ensuremath{\ulcorner}}
\providecommand{\disjointsum}{\mathbin{\oplus}}
\providecommand{\cardeq}[2]{#1 = #2}
\providecommand{\xred}{\mathrel{\rightarrow}}
% --- sixth compatibility block: names reported by real compile logs ---
\providecommand{\ORefut}{\operatorname{Refut}}
\providecommand{\ORProv}{\operatorname{Prov}}
\providecommand{\cutrank}{\operatorname{cr}}
\providecommand{\Refut}{\operatorname{Refut}}
\providecommand{\DischargeRule}{\mathrm{Dis}}
\providecommand{\Contraction}{\mathrm{C}}
\providecommand{\bottomAlignProof}{\relax}
\providecommand{\FalseCl}{\bot}
\providecommand{\red}{\mathrel{\rightarrow}}
\providecommand{\varAssign}{\sigma}
\providecommand{\defiff}{\mathrel{\leftrightarrow}}
\providecommand{\FalseInt}{\bot}
\providecommand{\fn}{\operatorname{fn}}
\providecommand{\dotted}{\cdot}

% --- tag-macro semantic fix (2026-10-13, final): \iftag CHOOSES a branch and the OLP proof body is
% always the SECOND group (\iftag{tag}{body}, {tag}{body}{c}, {tag}{body}{c}{d}); a definition that
% prints the tag literal drops the proof body. Bodies contain blank lines, hence \long.
\makeatletter
\providecommand{\iftag}{}
\long\def\iftag#1#2{\@ifnextchar\bgroup{\iftag@three{#2}}{\iftag@stop{#2}}}
\long\def\iftag@three#1#2{\iftag@stop{#1}\@ifnextchar\bgroup{\iftag@eat}{\relax}}
\long\def\iftag@eat#1{}
\long\def\iftag@stop#1{#1}
\long\def\oliflabeldef#1#2{\@ifnextchar\bgroup{\@gobble}{\iftag@stop{#2}}}
\providecommand{\tagrefs}[1]{}
\providecommand{\settag}[1]{}
\makeatother
\begin{document}

% ==== OLP file: content/normal-modal-logic/completeness/canonical-models.tex ====
% Part: normal-modal-logic
% Chapter: completeness
% Section: canonical-models



\olfileid{nml}{com}{cmd}

\olsection{Canonical Models}

The \emph{canonical model} for a modal system~$\Sigma$ is a specific
model~$\mModel{M^\Sigma}$ in which the worlds are all complete
$\Sigma$-consistent sets. Its accessibility relation~$R^\Sigma$ and
valuation $V^\Sigma$ are defined so as to guarantee that the
\olup{formula}s true at a world~$\Delta$ are exactly the \olup{formula}s
making up~$\Delta$.

\begin{defn}
  Let $\Sigma$ be a normal modal logic. The \emph{canonical model} for
  $\Sigma$ is $\mModel{M}^\Sigma = \tuple{W^\Sigma, R^\Sigma,
    V^\Sigma}$, where:
  \begin{enumerate}
  \item $W^\Sigma = \Setabs{ \Delta }{\Delta \text{
      is complete $\Sigma$-consistent} }$.
  \item $R^\Sigma \Delta\Delta'$ holds if and only if
    \iftag{prvBox}{$\Box^{-1}\Delta \subseteq
      \Delta'$}{$\Diamond\Delta' \subseteq \Delta$}.
  \item $V^\Sigma(p) = \Setabs{\Delta}{p \in \Delta}$.
  \end{enumerate}
\end{defn}



% ==== OLP file: content/normal-modal-logic/completeness/complete-consistent-sets.tex ====
% Part: normal-modal-logic
% Chapter: completeness
% Section: complete-consistent-sets



\olfileid{nml}{com}{ccs}

\olsection{Complete $\Sigma$-Consistent Sets}

Suppose $\Sigma$ is a set of modal \olup{formula}s---think of them as the
axioms or defining principles of a normal modal logic. A set $\Gamma$
is $\Sigma$-consistent iff $\Gamma \Proves/[\Sigma] \lfalse$, i.e., if
there is no \olup{derivation} of~$A_1 \lif (A_2 \lif \cdots (A_n \lif
\lfalse)\dots)$ from $\Sigma$, where each $A_i \in \Gamma$. We will
construct a ``canonical'' model in which each world is taken to be a
special kind of $\Sigma$-consistent set: one which is not just
$\Sigma$-consistent, but maximally so, in the sense that it settles
the truth value of every modal \olup{formula}: for every $A$, either $A
\in \Gamma$ or $\lnot A \in \Gamma$:

\begin{defn}
  A set $\Gamma$ is \emph{complete $\Sigma$-consistent} if and
  only if it is $\Sigma$-consistent and for every~$A$, either
  $A \in \Gamma$ or $\lnot A \in \Gamma$.
\end{defn}

Complete $\Sigma$-consistent sets $\Gamma$ have a number of useful
properties. For one, they are deductively closed, i.e., if $\Gamma
\Proves[\Sigma] A$ then $A \in \Gamma$. This means in particular
that every instance of \olup{formula}~$A \in \Sigma$ is also $\in
\Gamma$. Moreover, membership in $\Gamma$ mirrors the truth conditions
for the propositional connectives. This will be important when we
define the ``canonical model.''

\begin{prop}\ollabel{prop:ccs-properties}
  Suppose $\Gamma$ is complete $\Sigma$-consistent. Then:
  \begin{enumerate}
  \item \ollabel{prop:ccs-closed}%
    $\Gamma$ is deductively closed in~$\Sigma$.
  \item \ollabel{prop:ccs-sigma}%
    $\Sigma \subseteq \Gamma$.
  \tagitem{prvFalse}{\ollabel{prop:ccs-lfalse}%
    $\lfalse \notin \Gamma$}{}
  \tagitem{prvTrue}{\ollabel{prop:ccs-ltrue}%
    $\ltrue \in \Gamma$}{}
  \item \ollabel{prop:ccs-lnot}%
    $\lnot A \in \Gamma$ if and only if $A \notin
    \Gamma$.
  \tagitem{prvAnd}{\ollabel{prop:ccs-land}%
    $A \land B \in \Gamma$ iff $A \in \Gamma$ and $B \in \Gamma$}{}
  \tagitem{prvOr}{\ollabel{prop:ccs-lor}%
    $A \lor B \in \Gamma$ iff $A \in \Gamma$ or $B \in \Gamma$}{}
  \tagitem{prvIf}{\ollabel{prop:ccs-lif}%
    $A \lif B \in \Gamma$ iff $A \notin \Gamma$ or $B \in \Gamma$}{}
  \tagitem{prvIff}{\ollabel{prop:ccs-liff}%
    $A \liff B \in \Gamma$ iff either $A \in \Gamma$ and $B \in
    \Gamma$, or $A \notin \Gamma$ and $B \notin \Gamma$}{}
  \end{enumerate}
\end{prop}

\begin{proof}
  \begin{enumerate}
  \item Suppose $\Gamma \Proves[\Sigma] A$ but $A \notin
    \Gamma$. Then since $\Gamma$ is complete $\Sigma$-consistent,
    $\lnot A \in \Gamma$. This would make $\Gamma$ inconsistent, since
    $A, \lnot A \Proves[\Sigma] \lfalse$.
    
  \item If $A \in \Sigma$ then $\Gamma \Proves[\Sigma] A$, and $A
    \in \Gamma$ by deductive closure, i.e.,
    case~\olref{prop:ccs-closed}.
    
  \tagitem{prvFalse}{If $\lfalse \in \Gamma$, then $\Gamma
    \Proves[\Sigma] \lfalse$, so $\Gamma$ would be
    $\Sigma$-inconsistent.}{}
  
  \tagitem{prvTrue}{$\Gamma \Proves[\Sigma] \ltrue$, so $\ltrue \in
    \Gamma$ by deductive closure, i.e.,
    case~\olref{prop:ccs-closed}.}{}

\item If $\lnot A \in \Gamma$, then by consistency $A \notin
    \Gamma$; and if $A \notin \Gamma$ then $A \in \Gamma$ since
    $\Gamma$ is complete $\Sigma$-consistent.
    
  \tagitem{prvAnd}{\iftag{probAnd}{Exercise.}{Suppose $A \land B \in
      \Gamma$. Since $(A \land B) \lif A$ is a tautological
      instance, $A \in \Gamma$ by deductive closure, i.e.,
      case~\olref{prop:ccs-closed}. Similarly for $B \in \Gamma$. On
      the other hand, suppose both $A \in \Gamma$ and $B \in
      \Gamma$. Then deductive closure implies $(A \land B) \in
      \Gamma$, since $A \lif (B \lif (A \land B))$ is a
      tautological instance.}}{}

  \tagitem{prvOr}{\iftag{probOr}{Exercise.}{Suppose $A \lor B \in
      \Gamma$, and $A \notin \Gamma$ and $B \notin \Gamma$. Since
      $\Gamma$ is complete $\Sigma$-consistent, $\lnot A \in \Gamma$
      and $\lnot B \in \Gamma$. Then $\lnot(A \lor B) \in \Gamma$
      since $\lnot A \lif (\lnot B \lif \lnot (A \lor B))$ is a
      tautological instance. This would mean that $\Gamma$ is
      $\Sigma$-inconsistent, a contradiction.}}{}
  
  \tagitem{prvIf}{\iftag{probIf}{Exercise.}{Suppose $A \lif B \in
      \Gamma$ and $A \in \Gamma$; then $\Gamma \Proves[\Sigma] B$,
      whence $B \in \Gamma$ by deductive closure. Conversely, if $A
      \lif B \notin \Gamma$ then since $\Gamma$ is complete
      $\Sigma$-consistent, $\lnot (A \lif B) \in \Gamma$. Since
      $\lnot(A \lif B) \lif A$ is a tautological instance, $A \in
      \Gamma$ by deductive closure. Since $\lnot(A \lif B) \lif
      \lnot B$ is a tautological instance, $\lnot B \in
      \Gamma$. Then $B \notin \Gamma$ since $\Gamma$ is
      $\Sigma$-consistent.}}{}
  
  \tagitem{prvIff}{\iftag{probIff}{Exercise.}{Suppose $A \liff B \in
      \Gamma$. If $A \in \Gamma$, then $B \in \Gamma$, since $(A
      \liff B) \lif (A \lif B)$ is a tautological
      instance. Similarly, if $B \in \Gamma$, then $A \in \Gamma$.
      So either both $A \in \Gamma$ and $B \in \Gamma$, or neither
      $A \in \Gamma$ nor $B \in \Gamma$.
      
      Conversely, suppose $A \lif B \notin \Gamma$. Since $\Gamma$
      is complete $\Sigma$-consistent, $\lnot (A \liff B) \in
      \Gamma$. Since $\lnot(A \liff B) \lif (A \lif \lnot B)$ is a
      tautological instance, if $A \in \Gamma$ then $\lnot B \in
      \Gamma$, and since $\Gamma$ is $\Sigma$-consistent, $B \notin
      \Gamma$. Similarly, if $B \in \Gamma$ then $A \notin \Gamma$.
      So neither $A \in \Gamma$ and $B \in \Gamma$, nor $A \notin
      \Gamma$ and $B \notin \Gamma$.}}{}
  \end{enumerate}
\end{proof}

\begin{probtag}{probAnd,probOr,probIf,probIff}
  Complete the proof of \olref[nml][com][ccs]{prop:ccs-properties}.
\end{probtag}


% ==== OLP file: content/normal-modal-logic/completeness/completeness-K.tex ====
% Part: normal-modal-logic
% Chapter: completeness
% Section: completeness-K



\olfileid{nml}{com}{cmk}

\olsection{Determination and Completeness for  \Log{K}}

We are now prepared to use the canonical model to establish
completeness. Completeness follows from the fact that the \olup{formula}s
true in the canonical model for~$\Sigma$ are exactly the
$\Sigma$-\olup{derivable} ones. Models with this property are said to
\emph{determine}~$\Sigma$.

\begin{defn}
  A model $\mModel{M}$ \emph{determines} a normal modal
  logic~$\Sigma$ precisely when $\mSat{M}{A}$ if and only if $\Sigma
  \Proves A$, for all \olup{formula}s~$A$.
\end{defn}

\begin{thm}[Determination]\ollabel{thm:determination}
   $\mSat{M^\Sigma}{A}$ if and only if $\Sigma \Proves A$.
\end{thm}

\begin{proof}
  If $\mSat{M^\Sigma}{A}$, then for every complete
  $\Sigma$-consistent $\Delta$, we have
  $\mSat{M^\Sigma}{A}[\Delta]$. Hence, by the Truth Lemma, $A \in
  \Delta$ for every complete $\Sigma$-consistent $\Delta$, whence by
  \olref[lin]{cor:provability-characterization} (with $\Gamma =
  \emptyset$), $\Sigma \Proves A$.

  Conversely, if $\Sigma \Proves A$ then by
  \olref[ccs]{prop:ccs-properties}\olref[ccs]{prop:ccs-closed}, every
  complete $\Sigma$-consistent $\Delta$ contains~$A$, and hence by
  the Truth Lemma, $\mSat{M^\Sigma}{A}[\Delta]$ for every~$\Delta \in
  W^\Sigma$, i.e., $\mSat{M^\Sigma}{A}$.
\end{proof}

Since the canonical model for \Log{K} determines~\Log{K}, we
immediately have completeness of~\Log{K} as a corollary:

\begin{cor}\ollabel{cor:Kcomplete}
  The basic modal logic \Log{K} is complete with respect to the
  class of all models, i.e., if $\Entails A$ then $\Log{K}
  \Proves A$.
\end{cor}

\begin{proof}
  Contrapositively, if $\Log{K} \Proves/ A$ then by Determination
  $\mSat/{M^{\Log{K}}}{A}$ and hence $A$ is not valid.
\end{proof}

For the general case of completeness of a system $\Sigma$ with respect
to a class of models, e.g., of \Log{KTB4} with respect to the class of
reflexive, symmetric, transitive models, determination alone is not
enough.  We must also show that the canonical model for the system
$\Sigma$ is a member of the class, which does not follow obviously
from the canonical model construction---nor is it always true!


% ==== OLP file: content/normal-modal-logic/completeness/completeness.tex ====
% Part: normal-modal-logic
% Chapter: completeness



\olchapter{nml}{com}{Completeness and Canonical Models}

\olimport{introduction}
\olimport{complete-consistent-sets}
\olimport{lindenbaums-lemma}
\olimport{modalities-ccs}
\olimport{canonical-models}
\olimport{truth-lemma}
\olimport{completeness-K}
\olimport{frame-completeness}

\OLEndChapterHook



% ==== OLP file: content/normal-modal-logic/completeness/frame-completeness.tex ====
% Part: normal-modal-logic
% Chapter: completeness
% Section: frame-completeness



\olfileid{nml}{com}{fra}

\olsection{Frame Completeness}

The completeness theorem for \Log{K} can be extended to other modal
systems, once we show that the canonical model for a given logic has
the corresponding frame property.

\begin{thm}\ollabel{thm:completeframeprops}
  If a normal modal logic $\Sigma$ contains one of the \olup{formula}s on the
  left-hand side of \olref{tab:correspondencetable},
  then the canonical model for~$\Sigma$ has the corresponding property
  on the right-hand side.
\end{thm}

\begin{table}[htp]
  \centering
    \begin{tabular}{| l || l |}
      \hline
      \emph{If $\Sigma$ contains \dots} & \emph{\dots the canonical
        model for $\Sigma$ is:} \\
      \hline \hline
      \Ax{D}:\;\; $\Box A \lif \Diamond A$ & \quad serial; \\
      \hline
      \Ax{T}:\;\; $\Box A \lif A$ &\quad  reflexive;\\
      \hline
      \Ax{B}: \;\;$A \lif \Box\Diamond A$ &\quad  symmetric; \\
      \hline
      \Ax{4}: \;\; $\Box A \lif \Box\Box A$ & \quad transitive; \\
      \hline 
      \Ax{5}: \;\; $\Diamond A \lif \Box\Diamond A$& \quad euclidean.\\
      \hline
    \end{tabular}
    \caption{Basic correspondence facts.}
    \ollabel{tab:correspondencetable}
\end{table}

\begin{proof}
  We take each of these up in turn.

  Suppose $\Sigma$ contains \Ax{D}, and let $\Delta \in W^\Sigma$; we
  need to show that there is a $\Delta'$ such that $R^\Sigma
  \Delta\Delta'$. It suffices to show that $\Box^{-1}\Delta$ is
  $\Sigma$-consistent, for then by Lindenbaum's Lemma, there is a
  complete $\Sigma$-consistent set $\Delta' \supseteq
  \Box^{-1}\Delta$, and by definition of $R^\Sigma$
  \iftag{prvBox}{}{and \olref[mod]{lem:box-iff-diamond} }we have
  $R^\Sigma \Delta\Delta'$. So, suppose for contradiction that
  $\Box^{-1}\Delta$ is \emph{not} $\Sigma$-consistent, i.e.,
  $\Box^{-1}\Delta \Proves[\Sigma] \lfalse$. By \olref[mod]{lem:box2},
  $\Delta \Proves[\Sigma] \Box\lfalse$, and since $\Sigma$ contains
  \Ax{D}, also $\Delta \Proves[\Sigma] \Diamond\lfalse$. But $\Sigma$
  is normal, so $\Sigma \Proves \lnot\Diamond \lfalse$
  (\olref[prf][nor]{prop:notDiamondBot}), whence also $\Delta
  \Proves[\Sigma] \lnot\Diamond \lfalse$, against the consistency
  of~$\Delta$.

  Now suppose $\Sigma$ contains \Ax{T}, and let $\Delta \in
  W^\Sigma$. We want to show $R^\Sigma \Delta\Delta$,
  i.e.,\iftag{prvBox}{}{ by \olref[mod]{lem:box-iff-diamond},}
  $\Box^{-1}\Delta \subseteq \Delta$. But if $\Box A \in \Delta$ then
  by \Ax{T} also $A \in \Delta$, as desired.

  Now suppose $\Sigma$ contains \Ax{B}, and suppose $R^\Sigma
  \Delta\Delta'$ for $\Delta$, $\Delta' \in W^\Sigma$. We need to show
  that $R^\Sigma \Delta'\Delta$, i.e.,\iftag{prvBox}{
    $\Box^{-1}\Delta' \subseteq \Delta$. By
    \olref[mod]{lem:box-iff-diamond}, this is equivalent to}{}
  $\Diamond\Delta \subseteq \Delta'$. So suppose $A \in \Delta$. By
  \Ax{B}, also $\Box\Diamond A \in \Delta$. By the hypothesis that
  $R^\Sigma \Delta\Delta'$\iftag{prvBox}{}{ and
    \olref[mod]{lem:box-iff-diamond}}, we have that $\Box^{-1}\Delta
  \subseteq \Delta'$, and hence $\Diamond A \in \Delta'$, as required.

  Now suppose $\Sigma$ contains \Ax{4}, and suppose $R^\Sigma
  \Delta_1\Delta_2$ and $R^\Sigma \Delta_2\Delta_3$. We need to show
  $R^\Sigma \Delta_1\Delta_3$. From the hypothesis\iftag{prvBox}{}{ and
    \olref[mod]{lem:box-iff-diamond}} we have both $\Box^{-1}\Delta_1
  \subseteq \Delta_2$ and $\Box^{-1}\Delta_2 \subseteq \Delta_3$. In
  order to show $R^\Sigma \Delta_1\Delta_3$ it suffices to show
  $\Box^{-1}\Delta_1 \subseteq \Delta_3$. So let $B \in
  \Box^{-1}\Delta_1$, i.e., $\Box B \in \Delta_1$. By \Ax{4}, also
  $\Box\Box B \in \Delta_1$ and by hypothesis we get, first, that
  $\Box B \in \Delta_2$ and, second, that $B \in \Delta_3$, as
  desired.

  Now suppose $\Sigma$ contains \Ax{5}, suppose $R^\Sigma
  \Delta_1\Delta_2$ and $R^\Sigma \Delta_1\Delta_3$. We need to show
  $R^\Sigma \Delta_2\Delta_3$. The first hypothesis gives
  $\Box^{-1}\Delta_1 \subseteq \Delta_2$\iftag{prvBox}{}{ by
    \olref[mod]{lem:box-iff-diamond}}, and the second hypothesis is
  equivalent to $\Diamond\Delta_3 \subseteq \Delta_1$\iftag{prvBox}{,
    by \olref[mod]{lem:box-iff-diamond}}{}.  To show $R^\Sigma
  \Delta_2\Delta_3$\iftag{prvBox}{, by
    \olref[mod]{lem:box-iff-diamond}, it suffices to}{ we have to}
  show $\Diamond\Delta_3 \subseteq \Delta_2$. So let $\Diamond A \in
  \Diamond\Delta_3$, i.e., $A \in \Delta_3$. By the second hypothesis
  $\Diamond A \in \Delta_1$ and by \Ax{5}, $\Box\Diamond A \in
  \Delta_1$ as well. But now the first hypothesis gives $\Diamond A \in
  \Delta_2$, as desired.
  \end{proof}

As a corollary we obtain completeness results for a number of
systems. For instance, we know that $\Log{S5} = \Log{KT5} =
\Log{KTB4}$ is complete with respect to the class of all reflexive
euclidean models, which is the same as the class of all reflexive,
symmetric and transitive models.

\begin{thm}\ollabel{thm:generaldet}
  Let $\mClass{C}_\Ax{D}$, $\mClass{C}_\Ax{T}$,
  $\mClass{C}_\Ax{B}$, $\mClass{C}_\Ax{4}$, and
  $\mClass{C}_\Ax{5}$ be the class of all serial, reflexive,
  symmetric, transitive, and euclidean models (respectively). Then for
  any schemas $A_1$, \dots, $A_n$ among \Ax{D},
  \Ax{T}, \Ax{B}, \Ax{4}, and \Ax{5}, the system
  $\Log{K}A_1 \dots A_n$ is determined by the
  class of models $\mClass{C} = \mClass{C}_{A_1} \cap \dots
  \cap \mClass{C}_{A_n}$. 
\end{thm}

\begin{prop}
  Let $\Sigma$ be a normal modal logic; then:
  \begin{enumerate}
  \item\ollabel{prop:anotherfive-a}%
    If $\Sigma$ contains the schema $\Diamond A \lif \Box 
    A$ then the canonical model for $\Sigma$ is partially functional. 
  \item If $\Sigma$ contains the schema $\Diamond A \liff \Box 
    A$ then the canonical model for $\Sigma$ is functional. 
  \item If $\Sigma$ contains the schema $\Box\Box A \lif \Box 
    A$ then the canonical model for $\Sigma$ is weakly dense. 
  \end{enumerate}
(see \olref[frd][acc]{tab:anotherfive} for definitions of these frame
  properties).
\end{prop}

\begin{proof}
  \begin{enumerate}
  \item Suppose that $\Sigma$ contains the schema $\Diamond A \lif
    \Box A$, to show that $R^\Sigma$ is partially functional we need
    to prove that for any $\Delta_1$, $\Delta_2$, $\Delta_3 \in
    W^\Sigma$, if $R^\Sigma \Delta_1\Delta_2$ and $R^\Sigma
    \Delta_1\Delta_3$ then $\Delta_2=\Delta_3$. Since $R^\Sigma
    \Delta_1\Delta_2$ we have $\Box^{-1}\Delta_1 \subseteq \Delta_2$
    and since $R^\Sigma \Delta_1\Delta_3$ also $\Box^{-1}\Delta_1
    \subseteq \Delta_3$\iftag{prvBox}{}{, both by
      \olref[mod]{lem:box-iff-diamond}}. The identity
    $\Delta_2=\Delta_3$ will follow if we can establish the two
    inclusions $\Delta_2 \subseteq \Delta_3$ and $\Delta_3 \subseteq
    \Delta_2$. For the first inclusion, let $A \in \Delta_2$; then
    $\Diamond A \in \Delta_1$, and by the schema and deductive closure
    of $\Delta_1$ also $\Box A \in \Delta_1$, whence by the hypothesis
    that $R^\Sigma \Delta_1\Delta_3$, $A \in \Delta_3$. The second
    inclusion is similar.

  \item This follows immediately from part~\olref{prop:anotherfive-a}
    and the seriality proof in \olref{thm:completeframeprops}.

  \item Suppose $\Sigma$ contains the schema $\Box\Box A \lif \Box A$
    and to show that $R^\Sigma$ is weakly dense, let $R^\Sigma
    \Delta_1\Delta_2$. We need to show that there is a complete
    $\Sigma$-consistent set $\Delta_3$ such that $R^\Sigma
    \Delta_1\Delta_3$ and $R^\Sigma \Delta_3\Delta_2$. Let:
    \[
    \Gamma = \Box^{-1}\Delta_1 \cup \Diamond\Delta_2.
    \]
    It suffices to show that $\Gamma$ is $\Sigma$-consistent, for then
    by Lindenbaum's Lemma it can be extended to a complete
    $\Sigma$-consistent set~$\Delta_3$ such that $\Box^{-1}\Delta_1
    \subseteq \Delta_3$ and $\Diamond\Delta_2 \subseteq \Delta_3$,
    i.e., $R^\Sigma \Delta_1\Delta_3$\iftag{prvBox}{}{ (by
      \olref[mod]{lem:box-iff-diamond})} and $R^\Sigma
    \Delta_3\Delta_2$\iftag{prvBox}{ (by
      \olref[mod]{lem:box-iff-diamond})}{}.

    Suppose for contradiction that $\Gamma$ is not consistent. Then
    there are \olup{formula}s $\Box A_1$, \dots, $\Box A_n \in \Delta_1$
    and $B_1$, \dots,~$B_m \in \Delta_2$ such that \[A_1, \dots,
    A_n, \Diamond B_1, \dots, \Diamond B_m \Proves[\Sigma]
    \lfalse.\] Since $\Diamond (B_1 \land \dots \land B_m) \to
    (\Diamond B_1 \land \dots \land \Diamond B_m)$ is \olup{derivable} in
    every normal modal logic, we argue as follows, contradicting the
    consistency of~$\Delta_2$:
    \begin{align*}
      A_1, \dots, A_n, & \Diamond B_1, \dots,
      \Diamond B_m \Proves[\Sigma] \lfalse \\
      A_1, \dots,A_n & \Proves[\Sigma] (\Diamond B_1 \land
      \dots \land \Diamond B_m) \lif \lfalse\\
      & \qquad\text{by the deduction theorem}\\
      & \qquad\text{\olref[prf][prp]{prop:derivabilityfacts}\olref[prf][prp]{prop:derivabilityfacts-deduction}, and \Taut} \\
      A_1, \dots,A_n & \Proves[\Sigma]
      \Diamond(B_1 \land
      \dots \land B_m) \lif \lfalse\\
      & \qquad \text{since $\Sigma$ is normal} \\
      A_1, \dots,A_n & \Proves[\Sigma]
      \lnot\Diamond (B_1 \land \dots \land B_m)\\
      & \qquad\text{by \PL} \\
      A_1, \dots,A_n & \Proves[\Sigma]
      \Box\lnot (B_1 \land \dots \land B_m)\\
      & \qquad\text{$\Box\lnot$ for $\lnot\Diamond$} \\
      \Box A_1, \dots,\Box A_n
      & \Proves[\Sigma] \Box\Box \lnot (B_1 \land \dots \land B_m)\\
      & \qquad\text{by \olref[mod]{lem:box1}} \\
      \Box A_1, \dots,\Box A_n
      & \Proves[\Sigma]
      \Box\lnot (B_1 \land \dots \land B_m)\\
      &\qquad\text{by schema $\Box\Box A \lif \Box A$} \\
      \Delta_1 & \Proves[\Sigma]
      \Box\lnot (B_1 \land \dots \land B_m)\\
      &\qquad\text{by monotonicity, \olref[prf][prp]{prop:derivabilityfacts}\olref[prf][prp]{prop:derivabilityfacts-monotonicity}} \\
      & \Box\lnot (B_1 \land \dots \land B_m) \in \Delta_1\\
      &\qquad\text{by deductive closure}; \\
      & \lnot (B_1 \land \dots \land B_m) \in \Delta_2\\
      &\qquad \text{since }
      R^\Sigma \Delta_1\Delta_2. 
    \end{align*}
  \end{enumerate}
\end{proof}

On the strength of these examples, one might think that every
system~$\Sigma$ of modal logic is \emph{complete}, in the sense that
it proves every formula which is valid in every frame in which every
theorem of $\Sigma$ is valid. Unfortunately, there are many systems
that are not complete in this sense.


% ==== OLP file: content/normal-modal-logic/completeness/introduction.tex ====
% Part: normal-modal-logic
% Chapter: completeness
% Section: introduction



\olfileid{nml}{com}{int}

\olsection{Introduction}

If $\Sigma$ is a modal system, then the soundness theorem establishes
that if $\Sigma \Proves A$, then $A$ is valid in any class
$\mClass{C}$ of models in which all instances of all \olup{formula}s
in~$\Sigma$ are valid. In particular that means that if $\Log{K}
\Proves A$ then $A$ is true in all models; if $\Log{KT} \Proves A$
then $A$ is true in all reflexive models; if $\Log{KD} \Proves A$
then $A$ is true in all serial models, etc.

Completeness is the converse of soundness: that \Log{K} is complete
means that if \olup{formula}~$A$ is valid, $\Proves A$, for instance.
Proving completeness is a lot harder to do than proving soundness.  It
is useful, first, to consider the contrapositive: \Log{K} is complete
iff whenever $\Proves/ A$, there is a countermodel, i.e., a
model~$\mModel{M}$ such that $\mSat/{M}{A}$. Equivalently (negating
$A$), we could prove that whenever $\Proves/ \lnot A$, there is a
model of~$A$.  In the construction of such a model, we can use
information contained in $A$.  When we find models for specific
\olup{formula}s we often do the same: e.g., if we want to find a
countermodel to $p \lif \Box q$, we know that it has to contain a
world where $p$ is true and $\Box q$ is false. And a world where $\Box
q$ is false means there has to be a world accessible from it where $q$
is false. And that's all we need to know: which worlds make the
\olup{propositional variable}s true, and which worlds are accessible from
which worlds.

In the case of proving completeness, however, we don't have a specific
\olup{formula}~$A$ for which we are constructing a model. We want to
establish that a model exists for every $A$ such that
$\Proves/[\Sigma] \lnot A$. This is a minimal requirement, since
\emph{if} $\Proves[\Sigma] \lnot A$, by soundness, there is no model
for~$A$ (in which $\Sigma$ is true).  Now note that $\Proves/[\Sigma]
\lnot A$ iff $A$ is $\Sigma$-consistent.  (Recall that $\Sigma
\Proves/[\Sigma] \lnot A$ and $A \Proves/[\Sigma] \lfalse$ are
equivalent.)  So our task is to construct a model for every
$\Sigma$-consistent \olup{formula}.

The trick we'll use is to find a $\Sigma$-consistent set of
\olup{formula}s that contains~$A$, but also other formulas which tell us
what the world that makes $A$ true has to look like. Such sets are
\emph{complete} $\Sigma$-consistent sets. It's not enough to
construct a model with a single world to make~$A$ true, it will have
to contain multiple worlds and an accessibility relation. The complete
$\Sigma$-consistent set containing~$A$ will also contain other
\olup{formula}s of the form $\Box B$ and~$\Diamond C$. In all
accessible worlds, $B$ has to be true; in at least one, $C$ has to
be true. In order to accomplish this, we'll simply take \emph{all}
possible complete $\Sigma$-consistent sets as the basis for the set of
worlds. A tricky part will be to figure out when a complete
$\Sigma$-consistent set should count as being accessible from another
in our model.

We'll show that in the model so defined, $A$ is true at a
world---which is also a complete $\Sigma$-consistent set---iff $A$ is
\olup{element} of that set.  If $A$ is $\Sigma$-consistent, it will be
\olup{element} of at least one complete $\Sigma$-consistent set (a fact
we'll prove), and so there will be a world where $A$~is true. So we
will have a single model where every $\Sigma$-consistent
\olup{formula}~$A$ is true at some world.  This single model is the
\emph{canonical} model for~$\Sigma$.


% ==== OLP file: content/normal-modal-logic/completeness/lindenbaums-lemma.tex ====
% Part: normal-modal-logic
% Chapter: completeness
% Section: lindenbaums-lemma



\olfileid{nml}{com}{lin}

\olsection{Lindenbaum's Lemma}

Lindenbaum's Lemma establishes that every $\Sigma$-consistent set of
\olup{formula}s is contained in at least one \emph{complete}
$\Sigma$-consistent set. Our construction of the canonical model will
show that for each complete $\Sigma$-consistent set~$\Delta$, there is
a world in the canonical model where all and only the \olup{formula}s
in~$\Delta$ are true. So Lindenbaum's Lemma guarantees that every
$\Sigma$-consistent set is true at some world in the canonical model.

\begin{thm}[Lindenbaum's Lemma]\ollabel{thm:lindenbaum}
  If $\Gamma$ is $\Sigma$-consistent then there is a complete
  $\Sigma$-consistent set $\Delta$ extending~$\Gamma$.
\end{thm}

\begin{proof}
Let $A_0$, $A_1$, \dots{} be an exhaustive listing of all formulas
of the language (repetitions are allowed). For instance, start by
listing $\Obj p_0$, and at each stage $n \ge 1$ list the finitely many
formulas of length~$n$ using only variables among $\Obj p_0$,
\dots,~$\Obj p_n$. We define sets of \olup{formula}s $\Delta_n$ by
induction on~$n$, and we then set $\Delta = \bigcup_n \Delta_n$. We
first put $\Delta_0 = \Gamma$. Supposing that $\Delta_n$ has been
defined, we define $\Delta_{n+1}$ by:
\[
\Delta_{n+1} =
\begin{cases}
  \Delta_n \cup \{A_n\}, & \text{if $\Delta_n \cup \{ A_n\}$
    is $\Sigma$-consistent;} \\
  \Delta_n \cup \{ \lnot A_n\}, & \text{otherwise.}
\end{cases}
\]
Now let $\Delta = \bigcup_{n=0}^\infty \Delta_n$.

We have to show that this definition actually yields a set $\Delta$
with the required properties, i.e., $\Gamma \subseteq \Delta$ and
$\Delta$ is complete $\Sigma$-consistent.

It's obvious that $\Gamma \subseteq \Delta$, since $\Delta_0 \subseteq
\Delta$ by construction, and $\Delta_0 = \Gamma$. In fact, $\Delta_n
\subseteq \Delta$ for all~$n$, since $\Delta$ is the union of
all~$\Delta_n$. (Since in each step of the construction, we add
\olup{formula} to the set already constructed, $\Delta_n \subseteq
\Delta_{n+1}$, so since $\subseteq$ is transitive, $\Delta_n \subseteq
\Delta_{m}$ whenever $n \le m$.)  At each stage of the construction, we
either add $A_n$ or $\lnot A_n$, and every \olup{formula} appears (at
least once) in the list of all~$A_n$. So, for every $A$ either $A
\in \Delta$ or $\lnot A \in \Delta$, so $\Delta$ is complete by
definition.

Finally, we have to show, that $\Delta$ is $\Sigma$-consistent.  To do
this, we show that (a) if $\Delta$ were $\Sigma$-inconsistent, then
some $\Delta_n$ would be $\Sigma$-inconsistent, and (b) all $\Delta_n$
are $\Sigma$-consistent.

So suppose $\Delta$ were $\Sigma$-inconsistent. Then $\Delta
\Proves[\Sigma] \lfalse$, i.e., there are $A_1$, \dots,~$A_k \in
\Delta$ such that $\Sigma \Proves A_1 \lif (A_2 \lif \cdots (A_k
\lif \lfalse)\dots)$. Since $\Delta = \bigcup_{n=0}^\infty \Delta_n$, each
$A_i \in \Delta_{n_i}$ for some~$n_i$. Let $n$ be the largest of
these. Since $n_i \le n$, $\Delta_{n_i} \subseteq \Delta_n$. So, all
$A_i$ are in some~$\Delta_n$. This would mean $\Delta_n
\Proves[\Sigma] \lfalse$, i.e., $\Delta_n$ is $\Sigma$-inconsistent.

To show that each $\Delta_n$ is $\Sigma$-consistent, we use a simple
induction on~$n$. $\Delta_0 = \Gamma$, and we assumed $\Gamma$ was
$\Sigma$-consistent. So the claim holds for $n = 0$. Now suppose it
holds for $n$, i.e., $\Delta_n$ is $\Sigma$-consistent. $\Delta_{n+1}$
is either $\Delta_n \cup \{A_n\}$ if that is $\Sigma$-consistent,
otherwise it is $\Delta_n \cup \{\lnot A_n\}$. In the first case,
$\Delta_{n+1}$ is clearly $\Sigma$-consistent. However, by
\olref[prf][con]{prop:consistencyfacts}\olref[prf][con]{prop:consistencyfacts-c},
either $\Delta_n \cup \{A_n\}$ or $\Delta_n \cup \{\lnot A_n\}$ is
consistent, so $\Delta_{n+1}$ is consistent in the other case as well.
\end{proof}

\begin{cor}\ollabel{cor:provability-characterization}
  $\Gamma \Proves[\Sigma] A$ if and only if $A \in \Delta$ for
  each  complete $\Sigma$-consistent set $\Delta$ extending $\Gamma$
  (including when $\Gamma = \emptyset$, in which case we get another
  characterization of the modal system $\Sigma$.)
\end{cor}

\begin{proof}
  Suppose $\Gamma \Proves[\Sigma] A$, and let $\Delta$ be any
  complete $\Sigma$-consistent set extending $\Gamma$. If $A
  \notin \Delta$ then by maximality $\lnot A \in \Delta$ and so
  $\Delta \Proves[\Sigma] A$ (by monotonicity) and $\Delta
  \Proves[\Sigma] \lnot A$ (by reflexivity), and so $\Delta$ is
  inconsistent. Conversely if $\Gamma \Proves/[\Sigma] A$, then
  $\Gamma \cup \{ \lnot A\}$ is $\Sigma$-consistent, and by
  Lindenbaum's Lemma there is a complete consistent set $\Delta$
  extending $\Gamma \cup \{ \lnot A \}$. By consistency, $A
  \notin \Delta$.
\end{proof}


% ==== OLP file: content/normal-modal-logic/completeness/modalities-ccs.tex ====
% Part: normal-modal-logic
% Chapter: completeness
% Section: modalities-ccs



\olfileid{nml}{com}{mod}

\olsection{Modalities and Complete Consistent Sets}

\begin{explain}
When we construct a model~$\mModel{M^\Sigma}$ whose set of worlds is
given by the complete $\Sigma$-consistent sets~$\Delta$ in some normal
modal logic~$\Sigma$, we will also need to define an accessibility
relation~$R^\Sigma$ between such ``worlds.'' We want it to be the case
that the accessibility relation (and the assignment $V^\Sigma)$ are
defined in such a way that $\mSat{M^\Sigma}{A}[\Delta]$ iff $A \in
\Delta$. How should we do this?

Once the accessibility relation is defined, the definition of truth at
a world ensures that \iftag{prvBox}{$\mSat{M^\Sigma}{\Box A}[\Delta]$
  iff $\mSat{M^\Sigma}{A}[\Delta']$ for all $\Delta'$ such that
  $R^\Sigma\Delta\Delta'$}{$\mSat{M^\Sigma}{\Diamond A}[\Delta]$ iff
  $\mSat{M^\Sigma}{A}[\Delta']$ for at least one $\Delta'$ such that
  $R^\Sigma\Delta\Delta'$}.  The proof that
$\mSat{M^\Sigma}{A}[\Delta]$ iff $A \in \Delta$ requires that this
is true in particular for \olup{formula}s starting with a modal operator,
i.e., \iftag{prvBox}{$\mSat{M^\Sigma}{\Box A}[\Delta]$ iff $\Box A
  \in \Delta$}{$\mSat{M^\Sigma}{\Diamond A}[\Delta]$ iff $\Diamond A
  \in \Delta$}. Combining this requirement with the definition of
truth at a world for \iftag{prvBox}{$\Box A$}{$\Diamond A$} yields:
\[
\iftag{prvBox}{%
  \Box A \in \Delta \text{ iff } A \in \Delta' \text{
    for all $\Delta'$ with } R^\Sigma\Delta\Delta'}{%
  \Diamond A \in \Delta \text{ iff } A \in \Delta' \text{ for at
    least one } \Delta' \text{ with } R^\Sigma\Delta\Delta'}
\]
\iftag{prvBox}{Consider the left-to-right direction: it says that if
  $\Box A \in \Delta$, then $A \in \Delta'$ for any $A$ and any
  $\Delta'$ with $R^\Sigma\Delta\Delta'$. If we stipulate that
  $R^\Sigma\Delta\Delta'$ iff $A \in \Delta'$ for all $\Box A \in
  \Delta$, then this holds. We can write the condition on the right of
  the ``iff'' more compactly as: $\Setabs{A}{\Box A \in \Delta}
  \subseteq \Delta'$.}{Consider the right-to-left direction: it says
  that if $A \in \Delta'$ then $\Diamond A \in \Delta$, for any $A$
  and $\Delta'$ with $R^\Sigma\Delta\Delta'$. If we stipulate that
  $R^\Sigma\Delta\Delta'$ holds whenever $\Diamond A \in \Delta$ for
  all $A \in \Delta'$, then this holds.  We can write the condition
  on the right of the ``iff'' more compactly as:
  $\Setabs{\Diamond A}{A \in \Delta'} \subseteq \Delta$.}

So the question is: does this definition of $R^\Sigma$ in fact
guarantee that \iftag{prvBox}{$\Box A \in \Delta$ iff
  $\mSat{M^\Sigma}{\Box A}[\Delta]$? \iftag{prvDiamond}{Does it also
    guarantee that }{}}{}\iftag{prvDiamond}{$\Diamond A \in \Delta$
  iff $\mSat{M^\Sigma}{\Diamond A}[\Delta]$?}{} The next few results
will establish this.
\end{explain}

\begin{defn}
  If $\Gamma$ is a set of \olup{formula}s, let
  \begin{align*}
    \Box\Gamma & = \Setabs{\Box B}{B \in \Gamma}\\
    \Diamond\Gamma & = \Setabs{\Diamond B}{B \in \Gamma}\\
    \intertext{and}
    \Box^{-1}\Gamma & = \Setabs{B}{\Box B \in \Gamma}\\
    \Diamond^{-1}\Gamma & = \Setabs{B}{\Diamond B \in \Gamma}\\
  \end{align*}
\end{defn}

In other words, $\Box\Gamma$ is $\Gamma$ with $\Box$ in front of every
\olup{formula} in~$\Gamma$; $\Box^{-1}\Gamma$ is all the $\Box$'ed
\olup{formula}s of~$\Gamma$ with the initial $\Box$'s removed.  This
definition is not terribly important on its own, but will simplify the
notation considerably.

Note that $\Box\Box^{-1}\Gamma \subseteq \Gamma$:
\[
\Box\Box^{-1}\Gamma = \Setabs{\Box B}{\Box B \in \Gamma}
\]
i.e., it's just the set of all those \olup{formula}s of~$\Gamma$ that
start with~$\Box$.

\begin{lem}\ollabel{lem:box1}
  If $\Gamma \Proves[\Sigma] A$ then $\Box\Gamma \Proves[\Sigma] \Box
  A$.
\end{lem}

\begin{proof}
  If $\Gamma \Proves[\Sigma] A$ then there are $B_1$, \dots, $B_k
  \in \Gamma$ such that $\Sigma \Proves B_1 \lif (B_2 \lif \cdots
  (B_n \lif A)\cdots)$. Since $\Sigma$ is normal, by rule \RK{},
  $\Sigma \Proves \Box B_1 \lif (\Box B_2 \lif \cdots (\Box B_n \lif
  \Box A)\cdots)$, where obviously $\Box B_1$, \dots, $\Box B_k \in
  \Box\Gamma$. Hence, by definition, $\Box\Gamma \Proves[\Sigma] \Box
  A$.
\end{proof}

\begin{lem}\ollabel{lem:box2}
  If $\Box^{-1} \Gamma \Proves[\Sigma] A$ then $\Gamma
  \Proves[\Sigma] \Box A$.
\end{lem}

\begin{proof}
  Suppose $\Box^{-1}\Gamma \Proves[\Sigma] A$; then by
  \olref{lem:box1}, $\Box\Box^{-1}\Gamma \Proves \Box A$. But since
  $\Box\Box^{-1}\Gamma \subseteq \Gamma$, also $\Gamma \Proves[\Sigma]
  \Box A$ by monotonicity.
\end{proof}

\iftag{prvBox}{%
% Box is primitive
  
\begin{prop}\ollabel{prop:box}
  If $\Gamma$ is complete $\Sigma$-consistent, then $\Box A \in
  \Gamma$ if and only if for every complete $\Sigma$-consistent
  $\Delta$ such that $\Box^{-1} \Gamma \subseteq \Delta$, it holds
  that $A \in \Delta$.
\end{prop}

\begin{proof}
  Suppose $\Gamma$ is complete $\Sigma$-consistent. The ``only if''
  direction is easy: Suppose $\Box A \in \Gamma$ and that
  $\Box^{-1}\Gamma \subseteq \Delta$. Since $\Box A \in \Gamma$, $A
  \in \Box^{-1}\Gamma \subseteq \Delta$, so $A \in \Delta$.

  For the ``if'' direction, we prove the contrapositive: Suppose
  $\Box A \notin \Gamma$. Since $\Gamma$ is complete
  $\Sigma$-consistent, it is deductively closed, and hence $\Gamma
  \Proves/[\Sigma] \Box A$. By \olref{lem:box2},
  $\Box^{-1}\Gamma \Proves/[\Sigma] A$. By
  \olref[prf][con]{prop:consistencyfacts}\olref[prf][con]{prop:consistencyfacts-b},
  $\Box^{-1}\Gamma \cup \{ \lnot A \}$ is $\Sigma$-consistent. By
  Lindenbaum's Lemma, there is a complete $\Sigma$-consistent
  set~$\Delta$ such that $\Box^{-1}\Gamma \cup \{ \lnot A \} \subseteq
  \Delta$. By consistency, $A \notin \Delta$.
\end{proof}
}{%
% Box is not primitive, only Diamond is

\begin{prop}\ollabel{prop:diamond}
  If $\Gamma$ is complete $\Sigma$-consistent, then $\Diamond A \in
  \Gamma$ if and only if for some complete $\Sigma$-consistent
  $\Delta$ such that $\Diamond\Delta \subseteq
  \Gamma$, it holds that $A \in \Delta$.
\end{prop}

\begin{proof}
  Suppose $\Gamma$ is complete $\Sigma$-consistent. The right-to-left
  part is easy: Let $\Delta$ be complete $\Sigma$-consistent such that
  $\Diamond\Delta \subseteq \Gamma$ and $A \in \Delta$.  Then
  $\Diamond A \in \Diamond\Delta \subseteq \Gamma$, i.e., $\Diamond A
  \in \Gamma$.

  For the left-to-right direction, assume $\Diamond A \in \Gamma$. We
  have to show that there is a complete $\Sigma$-consistent $\Delta$
  with $A \in \Delta$ and $\Diamond\Delta \subseteq \Gamma$.  Since
  $\Diamond A \in \Gamma$, and $\Gamma$ is $\Sigma$-consistent,
  $\Box\lnot A \notin \Gamma$. Since $\Gamma$ is complete
  $\Sigma$-consistent, it is deductively closed, so $\Gamma
  \Proves/[\Sigma] \Box\lnot A$.  By \olref{lem:box2},
  $\Box^{-1}\Gamma \Proves/[\Sigma] \lnot A$. By
  \olref[prf][con]{prop:consistencyfacts}\olref[prf][con]{prop:consistencyfacts-b},
  $\Box^{-1}\Gamma \cup \{ A \}$ is $\Sigma$-consistent. By
  Lindenbaum's Lemma, there is a complete $\Sigma$-consistent
  set~$\Delta$ such that $\Box^{-1}\Gamma \cup \{ A \} \subseteq
  \Delta$. Clearly, $A \in \Delta$.  To see that $\Diamond\Delta
  \subseteq \Gamma$, assume it weren't: For some $B \in \Delta$,
  $\Diamond B \notin \Gamma$. Then, since $\Gamma$ is complete
  $\Sigma$-consistent, $\Box\lnot B \in \Gamma$. But then also $\lnot
  B \in \Box^{-1}\Gamma \subseteq \Delta$, which would make $\Delta$
  inconsistent.
\end{proof}
}

\begin{lem}\ollabel{lem:box-iff-diamond}
  Suppose $\Gamma$ and $\Delta$ are complete
  $\Sigma$-consistent. Then $\Box^{-1}\Gamma \subseteq \Delta$ if and
  only if $\Diamond\Delta \subseteq \Gamma$.
\end{lem}

\begin{proof}
  ``Only if'' direction: Assume $\Box^{-1}\Gamma \subseteq \Delta$ and
  suppose $\Diamond A \in \Diamond\Delta$ (i.e., $A \in \Delta$).  In
  order to show $\Diamond A \in \Gamma$, it suffices to show
  $\Box\lnot A \notin \Gamma$, for then by maximality, $\lnot\Box\lnot
  A \in \Gamma$. Now, if $\Box\lnot A \in \Gamma$ then by hypothesis
  $\lnot A \in \Delta$, against the consistency of $\Delta$ (since $A
  \in \Delta$). Hence $\Box\lnot A \notin \Gamma$, as required.

  ``If'' direction: Assume $\Diamond\Delta
  \subseteq \Gamma$. We argue contrapositively: suppose $A
  \notin \Delta$ in order to show $\Box A \notin \Gamma$. If
  $A \notin \Delta$ then by maximality $\lnot A \in \Delta$
  and so by hypothesis $\Diamond\lnot A \in \Gamma$. But in a
  normal modal logic $\Diamond\lnot A$ is equivalent to
  $\lnot\Box A$, and if the latter is in $\Gamma$, by consistency
  $\Box A \notin\Gamma$, as required.
\end{proof}

\iftag{prvBox}{%
  \iftag{prvDiamond}{%
% Box and Diamond are both primitive; prove
% prop:diamond using lem:box-iff-diamond
  
\begin{prop}\ollabel{prop:diamond}
  If $\Gamma$ is complete $\Sigma$-consistent, then $\Diamond A \in
  \Gamma$ if and only if for some complete $\Sigma$-consistent
  $\Delta$ such that $\Diamond\Delta \subseteq
  \Gamma$, it holds that $A \in \Delta$.
\end{prop}

\begin{proof}
Suppose $\Gamma$ is complete $\Sigma$-consistent.  $\Diamond A \in
\Gamma$ iff $\lnot\Box\lnot A \in \Gamma$ by \Dual{} and closure.
$\lnot\Box\lnot A \in \Gamma$ iff $\Box\lnot A \notin \Gamma$ by
\olref[ccs]{prop:ccs-properties}\olref[ccs]{prop:ccs-lnot} since
$\Gamma$ is complete $\Sigma$-consistent.  By \olref{prop:box},
$\Box\lnot A \notin \Gamma$ iff, for some complete
$\Sigma$-consistent $\Delta$ with $\Box^{-1}\Gamma \subseteq \Delta$,
$\lnot A \notin \Delta$.  Now consider any such~$\Delta$. By
\olref{lem:box-iff-diamond}, $\Box^{-1}\Gamma \subseteq \Delta$ iff
$\Diamond\Delta \subseteq \Gamma$. Also, $\lnot A \notin \Delta$ iff
$A \in \Delta$ by
\olref[ccs]{prop:ccs-properties}\olref[ccs]{prop:ccs-lnot}.  So
$\Diamond A \in \Gamma$ iff, for some complete $\Sigma$-consistent
$\Delta$ with $\Diamond\Delta \subseteq \Gamma$, $A \in \Delta$.
\end{proof}

}{}}{}
  
\begin{prob}
  Show that if $\Gamma$ is complete $\Sigma$-consistent, then
  \iftag{prvBox}{$\Diamond A \in \Gamma$ if and only if there is a
    complete $\Sigma$-consistent $\Delta$ such that $\Box^{-1}\Gamma
    \subseteq \Delta$ and $A \in \Delta$.}{$\Box A \in
    \Gamma$ if and only if for every complete $\Sigma$-consistent
    $\Delta$ such that $\Box^{-1} \Gamma \subseteq \Delta$, it holds
    that $A \in \Delta$.}  \emph{Do this without using
    \olref[nml][com][mod]{lem:box-iff-diamond}}.
\end{prob}


% ==== OLP file: content/normal-modal-logic/completeness/truth-lemma.tex ====
% Part: normal-modal-logic
% Chapter: completeness
% Section: truth-lemma



\olfileid{nml}{com}{tru}

\olsection{The Truth Lemma}

The canonical model~$\mModel{M^\Sigma}$ is defined in such a way that
$\mSat{M^\Sigma}{A}[\Delta]$ iff $A \in \Delta$. For \olup{propositional
variable}s, the definition of $V^\Sigma$ yields this directly. We have
to verify that the equivalence holds for all \olup{formula}s, however. We
do this by induction. The inductive step involves proving the
equivalence for \olup{formula}s involving propositional operators (where
we have to use \olref[ccs]{prop:ccs-properties}) and the modal
operators (where we invoke the results of \olref[mod]{sec}).

\begin{prop}[Truth Lemma]
  \ollabel{prop:truthlemma}
  For every \olup{formula}~$A$, $\mSat{M^\Sigma}{A}[\Delta]$ if and only if
  $A \in \Delta$.
\end{prop}

\begin{proof}
  By induction on~$A$.
  \begin{enumerate}
    \tagitem{prvFalse}{\indcase{A}{\lfalse}{$\mSat/{M^\Sigma}{\lfalse}[\Delta]$
        by \olref[syn][trw]{defn:mmodels}, and $\lfalse \notin \Delta$
        by
        \olref[ccs]{prop:ccs-properties}\olref[ccs]{prop:ccs-lfalse}.}}{}

    \tagitem{prvTrue}{\indcase{A}{\ltrue}{$\mSat{M^\Sigma}{\ltrue}[\Delta]$
        by \olref[syn][trw]{defn:mmodels}, and $\ltrue \in \Delta$ by
        \olref[ccs]{prop:ccs-properties}\olref[ccs]{prop:ccs-ltrue}.}}{}

    \item \indcase{A}{p}{$\mSat{M^\Sigma}{p}[\Delta]$ iff $\Delta \in
      V^\Sigma(p)$ by \olref[syn][trw]{defn:mmodels}.  Also, $\Delta
      \in V^\Sigma(p)$ iff $p \in \Delta$ by definition
      of~$V^\Sigma$.}

    \tagitem{prvNot}{\indcase{A}{\lnot B}{%
        \iftag{probNot}{Exercise.}{$\mSat{M^\Sigma}{\lnot B}[\Delta]$
          iff $\mSat/{M^\Sigma}{B}[\Delta]$
          (\olref[syn][trw]{defn:mmodels}) iff $B \notin \Delta$ (by
          inductive hypothesis) iff $\lnot B \in \Delta$ (by
          \olref[ccs]{prop:ccs-properties}\olref[ccs]{prop:ccs-lnot}).}}}{}

    \tagitem{prvAnd}{\indcase{A}{B \land C}{%
        \iftag{probAnd}{Exercise.}{$\mSat{M^\Sigma}{B \land
            C}[\Delta]$ iff $\mSat{M^\Sigma}{B}[\Delta]$ and
          $\mSat{M^\Sigma}{C}[\Delta]$ (by
          \olref[syn][trw]{defn:mmodels}) iff $B \in \Delta$ and $C
          \in \Delta$ (by inductive hypothesis) iff $B \land C \in
          \Delta$ (by
          \olref[ccs]{prop:ccs-properties}\olref[ccs]{prop:ccs-land}).}}}{}

    \tagitem{prvOr}{\indcase{A}{B \lor C}{%
        \iftag{probOr}{Exercise.}{$\mSat{M^\Sigma}{B \lor
            C}[\Delta]$ iff $\mSat{M^\Sigma}{B}[\Delta]$ or
          $\mSat{M^\Sigma}{C}[\Delta]$ (by
          \olref[syn][trw]{defn:mmodels}) iff $B \in \Delta$ or $C
          \in \Delta$ (by inductive hypothesis) iff $B \lor C \in
          \Delta$ (by
          \olref[ccs]{prop:ccs-properties}\olref[ccs]{prop:ccs-lor}).}}}{}

    \tagitem{prvIf}{\indcase{A}{B \lif C}{%
        \iftag{probIf}{Exercise.}{$\mSat{M^\Sigma}{B \lif
            C}[\Delta]$ iff $\mSat/{M^\Sigma}{B}[\Delta]$ or
          $\mSat{M^\Sigma}{C}[\Delta]$ (by
          \olref[syn][trw]{defn:mmodels}) iff $B \notin \Delta$ or
          $C \in \Delta$ (by inductive hypothesis) iff $B \lif C
          \in \Delta$ (by
          \olref[ccs]{prop:ccs-properties}\olref[ccs]{prop:ccs-lif}).}}}{}

    \tagitem{prvIff}{\indcase{A}{B \liff C}{%
        \iftag{probIff}{Exercise.}{$\mSat{M^\Sigma}{B \liff
            C}[\Delta]$ iff either $\mSat{M^\Sigma}{B}[\Delta]$ and
          $\mSat{M^\Sigma}{C}[\Delta]$ or
          $\mSat/{M^\Sigma}{B}[\Delta]$ and
          $\mSat/{M^\Sigma}{C}[\Delta]$ (by
          \olref[syn][trw]{defn:mmodels}) iff either $B \in \Delta$
          and $C \in \Delta$ or $B \notin \Delta$ and $C \notin
          \Delta$ (by inductive hypothesis) iff $B \liff C \in
          \Delta$ (by
          \olref[ccs]{prop:ccs-properties}\olref[ccs]{prop:ccs-liff}).}}}{}

    \tagitem{prvBox}{\indcase{A}{\Box B}{%
        \iftag{probBox}{Exercise.}{First suppose that
          $\mSat{M^\Sigma}{\Box B}[\Delta]$. By
          \olref[syn][trw]{defn:mmodels}, for every $\Delta'$ such
          that $R^\Sigma \Delta\Delta'$,
          $\mSat{M^\Sigma}{B}[\Delta']$. By inductive hypothesis, for
          every $\Delta'$ such that $R^\Sigma \Delta\Delta'$, $B \in
          \Delta'$. By definition of $R^\Sigma$, for every $\Delta'$
          such that $\Box^{-1} \Delta \subseteq \Delta'$, $B \in
          \Delta'$. By \olref[mod]{prop:box}, $\Box B \in
          \Delta$.

          Now assume $\Box B \in \Delta$. Let $\Delta' \in W^\Sigma$
          be such that $R^\Sigma \Delta\Delta'$, i.e., $\Box^{-1}\Delta
          \subseteq \Delta'$. Since $\Box B \in \Delta$, $B \in
          \Box^{-1} \Delta$. Consequently, $B \in \Delta'$. By
          inductive hypothesis, $\mSat{M^\Sigma}{B}[\Delta']$. Since
          $\Delta'$ is arbitrary with $R^\Sigma \Delta\Delta'$, for
          all $\Delta' \in W^\Sigma$ such that $R^\Sigma
          \Delta\Delta'$, $\mSat{M^\Sigma}{B}[\Delta']$. By
          \olref[syn][trw]{defn:mmodels}, $\mSat{M^\Sigma}{\Box
            B}[\Delta]$.}}}{}

    \tagitem{prvDiamond}{\indcase{A}{\Diamond B}{%
        \iftag{probDiamond}{Exercise.}{First suppose that
          $\mSat{M^\Sigma}{\Diamond B}[\Delta]$.  By
          \olref[syn][trw]{defn:mmodels}, for some $\Delta'$ such that
          $R^\Sigma \Delta\Delta'$, $\mSat{M^\Sigma}{B}[\Delta']$.
          By inductive hypothesis, for some $\Delta'$ such that
          $R^\Sigma \Delta\Delta'$, $B \in \Delta'$. By definition of
          $R^\Sigma$,\iftag{prvBox}{ for some $\Delta'$ such that
            $\Box^{-1} \Delta \subseteq \Delta'$, $B \in \Delta'$.
            By \olref[mod]{prop:diamond},}{} for some $\Delta'$ such
          that $\Diamond \Delta' \subseteq \Delta$, $B \in
          \Delta'$. Since $B \in \Delta'$, $\Diamond B \in \Diamond
          \Delta'$, so $\Diamond B \in \Delta$.

          Now assume $\Diamond B \in \Delta$. By
          \olref[mod]{prop:diamond}, there is a complete
          $\Sigma$-consistent $\Delta' \in W^\Sigma$ such that
          $\Diamond \Delta' \subseteq \Delta$ and $B \in
          \Delta'$.\iftag{prvBox}{ By \olref[mod]{lem:box-iff-diamond},
            there is a $\Delta' \in W^\Sigma$ such that $\Box^{-1}
            \Delta \subseteq \Delta'$, and $B \in \Delta'$.}{} By
          definition of $R^\Sigma$, $R^\Sigma \Delta\Delta'$, so there
          is a $\Delta' \in W^\Sigma$ such that $R^\Sigma
          \Delta\Delta'$ and $B \in \Delta'$. By
          \olref[syn][trw]{defn:mmodels}, $\mSat{M^\Sigma}{\Diamond
            B}[\Delta]$.}}}{}
  \end{enumerate}
\end{proof}

\begin{probtag}{probFalse,probTrue,probNot,proband,probOr,probIf,probIff,probBox,probDiamond}
  Complete the proof of \olref[nml][com][tru]{prop:truthlemma}.
\end{probtag}


\end{document}
